\chapter{Picard 迭代}

\begin{solution}{picard-integral-equivalence}
若 \(y\in C^1\) 满足微分方程，微积分基本定理给
\[
 y(t)-y(t_0)=\int_{t_0}^{t}y'(s)ds
 =\int_{t_0}^{t}f(s,y(s))ds.
\]
反之，连续 \(f\) 与连续 \(y\) 使复合函数 \(s\mapsto f(s,y(s))\) 连续；
积分上限函数可导，导数为 \(f(t,y(t))\)，且令 \(t=t_0\) 得初值。
\end{solution}

\begin{solution}{picard-invariant-ball}
若 \(\|u-y_0\|_\infty\le b\)，则 \((t,u(t))\in Q\)，因而
\[
 \|(\Phi u)(t)-y_0\|
 \le\left|\int_{t_0}^{t}\|f(s,u(s))\|ds\right|
 \le M|t-t_0|\le Mh.
\]
所以 \(Mh\le b\) 保证 \(\Phi u\in X\)。
\end{solution}

\begin{solution}{picard-linear-equation}
从 \(u_0=y_0\) 开始，
\[
 u_1=y_0+(ay_0+b)t,\qquad
 u_2=y_0+(ay_0+b)\left(t+\frac{at^2}{2}\right),
\]
\[
 u_3=y_0+(ay_0+b)\left(t+\frac{at^2}{2}+\frac{a^2t^3}{6}\right).
\]
依此归纳，
\[
 u_n=y_0+(ay_0+b)\sum_{k=1}^n\frac{a^{k-1}t^k}{k!}.
\]
若 \(a\ne0\)，极限为
\((y_0+b/a)e^{at}-b/a\)；若 \(a=0\)，为 \(y_0+bt\)。
\end{solution}

\begin{solution}{picard-contraction-error}
压缩常数 \(q=Lh<1\) 时，
\[
 \|u_n-y\|\le\sum_{j=n}^\infty\|u_{j+1}-u_j\|
 \le\frac{q^n}{1-q}\|u_1-u_0\|.
\]
而 \(\|u_1-u_0\|\le Mh\)。阶乘估计从
\(\|u_1(t)-u_0(t)\|\le M|t-t_0|\) 开始。若第 \(n\) 步有逐点上界
\[
 \|u_{n+1}(t)-u_n(t)\|
 \le\frac{ML^n|t-t_0|^{n+1}}{(n+1)!},
\]
再积分一次即得下一步；对 \(|t-t_0|\le h\) 取上确界。
\end{solution}

\begin{solution}{picard-local-lipschitz-derivative}
对固定 \(t\) 和 \(y,z\) 的线段，有限维中值积分公式给
\[
 f(t,y)-f(t,z)=\int_0^1D_yf(t,z+s(y-z))(y-z)ds.
\]
连续导数在紧矩形上范数有界，记上界为 \(L\)，便得
\(\|f(t,y)-f(t,z)\|\le L\|y-z\|\)，且 \(L\) 与 \(t,y,z\) 无关。
\end{solution}

\begin{solution}{picard-nonunique-family}
\(t<c\) 时 \(y_c'=0=2\sqrt{|y_c|}\)；\(t>c\) 时
\(y_c'=2(t-c)=2\sqrt{(t-c)^2}\)。在 \(t=c\) 左右导数均为零，故
\(y_c\in C^1\)，方程也成立。所有 \(c\ge0\) 的解都取同一初值，却在
\(t>c\) 后不同，故唯一性失败。
\end{solution}

\begin{solution}{picard-gronwall-variable}
令 \(v=A+\int b u\)。当 \(A>0\) 时 \(u\le v\) 且
\((\log v)'=v'/v\le b\) 几乎处处，积分得
\(v(t)\le A\exp(\int b)\)。\(A=0\) 用 \(A+\varepsilon\) 后取极限。
对解差 \(u(t)=\|y(t)-z(t)\|\)，积分方程给
\[
 u(t)\le\|y_0-z_0\|+\int_{t_0}^{t}Lu(s)ds,
\]
应用引理即得所求。
\end{solution}

\begin{solution}{picard-perturbed-vector-field}
把差拆为
\[
 f(s,y(s))-g(s,z(s))
 =[f(s,y(s))-f(s,z(s))]+[f(s,z(s))-g(s,z(s))].
\]
因此
\[
 u(t)\le\|y_0-z_0\|+\eta|t-t_0|+
 L\int_{t_0}^{t}u(s)ds.
\]
在长度 \(T=|t-t_0|\) 的区间上把前两项统一界为
\(\|y_0-z_0\|+\eta T\)，Grönwall 给正文估计。
\end{solution}

\begin{solution}{picard-global-linear-growth}
由积分方程，
\[
 \|y(t)\|\le\|y_0\|+A|t-t_0|+
 B\int_{t_0}^{t}\|y(s)\|ds.
\]
Grönwall 给任意固定有限时间带上的共同有限上界。因此轨道不可能因
\(\|y(t)\|\to\infty\) 在有限端点终止。向量场定义在全空间，有限时间带内的
轨道闭包是紧集，局部存在长度可统一选择，最大解可继续延拓。
\end{solution}

\begin{solution}{picard-flow-semigroup}
记从初值 \(x\) 出发的唯一解为 \(y(t;x)\)，并令
\(\Phi_t(x)=y(t;x)\)。函数
\(s\mapsto y(s+t;x)\) 与 \(s\mapsto y(s;\Phi_t(x))\) 满足同一自治方程，
在 \(s=0\) 取同一初值。唯一性给
\[
 \Phi_s(\Phi_t(x))=\Phi_{s+t}(x).
\]
等式只在三项均有定义时成立。局部流的定义域通常是
\[
 \{(t,x):t\text{ 属于从 }x\text{ 出发的最大区间}\},
\]
而不是整个 \(\R\times\R^m\)。
\end{solution}

\begin{solution}{picard-parameter-differentiation}
设 \(y(t,\lambda)\) 解
\[
 y=y_0(\lambda)+\int_{t_0}^{t}f(s,y(s,\lambda),\lambda)ds.
\]
形式导数 \(z=\partial_\lambda y\) 应满足
\[
 z(t)=y_0'(\lambda)+\int_{t_0}^{t}
 \bigl[D_yf(s,y,\lambda)z(s)+\partial_\lambda f(s,y,\lambda)\bigr]ds.
\]
对差商与该线性积分方程之差使用一阶展开；连续可微性在紧管道上使余项除以
参数增量后一致趋零。差的范数满足
\(u(t)\le o(1)+L\int u\)，Grönwall 给 \(u\to0\) 一致。因此参数导数存在并
由变分方程唯一确定。
\end{solution}

\begin{solution}{picard-weighted-norm}
对 \(t\ge t_0\)，
\[
 e^{-\alpha(t-t_0)}\|(\Phi u)(t)-(\Phi v)(t)\|
 \le L\|u-v\|_\alpha e^{-\alpha(t-t_0)}
 \int_{t_0}^{t}e^{\alpha(s-t_0)}ds
 \le\frac L\alpha\|u-v\|_\alpha.
\]
\(\alpha>L\) 时压缩常数小于一。这个范数不能代替闭集不变估计，因为
\(f\) 的 Lipschitz 与有界条件仍只在预先选定的区域内已知。
\end{solution}
